\section{Discussion and conclusion}

Exchange retuning distinguishes two physically different origins of a
dark low-energy spectrum. For a local doublet, the first-order correction
has the same operator form as the projected dipole probe.
A single retuning cancels it throughout the retained space, suppressing
both ground-state and thermal weights to fourth order or smaller.
For the pentamer quartet, virtual excitations generate a composite vector
with different amplitudes in different bond-spin channels.
The positive minimum of its commutator moment forces a leading thermal
response at every positive scaled temperature on finite even rings.

The four-cell ring shows why a zero-temperature spectrum alone cannot
distinguish these mechanisms. Its leading quartet probe annihilates the
ground state after retuning, although the full operator remains dynamical.
Population of the first triplet restores lines at $2J_{\rm eff}$ and
$7J_{\rm eff}$ with leading intensity ratio $9:8$.
The twenty-spin microscopic spectrum approaches these line positions
and strengths in the weak-coupling limit.
Thus thermal scattering tests the compatibility of cancellation across
initial states, beyond the transition amplitudes sampled by the ground state.

The control-space result gives this distinction a direct physical meaning.
Independent first-order SU(2)-invariant changes within each cell adjust
local dipole coefficients while preserving the leading uniform chain.
Their nonuniform components add weight at other operator momenta, leaving
the optimal $Q=\pi$ residual unchanged.
The obstruction is therefore tied to the accessible operator structure
under these local controls. Altering the leading intercell interactions,
introducing anisotropy, or changing the local multiplet would change
that structure and poses a different control problem.

At finite coupling and momentum acceptance, the thermal signal competes
with two calculable backgrounds. A fourth-order ground-state contribution
dominates at sufficiently low temperature, while averaging around the
interference root restores ordinary form-factor weight.
Optimizing exchange over the momentum window removes its mean shift but
leaves its variance. These effects set the scales needed to interpret
a suppressed spectrum and separate thermal activation from residual
ground-state scattering.

The results establish a weak-coupling distinction at fixed finite size
between cancellation of a probe and cancellation for one initial state.
The effective uniform-trace bound further retains finite-frequency
quartet weight uniformly in even chain length, with the trace limit
taken first. Extending that bound to fixed scaled temperature in the
thermodynamic limit would connect the finite-chain mechanism to bulk
thermal scattering. Within the regime established here, the local
spin representation and accessible exchange controls determine whether
a retuned interference zero survives thermal occupation.
